% =============================================================================
%  Definiciones
% =============================================================================

\newdefinition{backend}{backend}{backends}%
  {Conjunto de componentes de un sistema responsables de la lógica de negocio,
   la persistencia de datos y la exposición de servicios, sin contacto directo
   con la interfaz de usuario}

\newdefinition{frontend}{frontend}{frontends}%
  {Capa del sistema responsable de la interacción con el usuario y de la
   presentación visual de la información}

\newdefinition{multitenant}{multi-tenant}{multi-tenant}%
  {Arquitectura que permite atender a múltiples organizaciones (tenants)
   desde una única instancia del sistema, manteniendo aislados sus datos y
   configuraciones}

\newdefinition{plugin}{plugin}{plugins}%
  {Componente intercambiable que extiende la funcionalidad del sistema
   cumpliendo una interfaz definida, sin requerir modificaciones en el resto
   del código}

\newdefinition{tenant}{tenant}{tenants}%
  {Cada una de las organizaciones que comparten una instancia del sistema en
   una arquitectura multi-tenant}

\newdefinition{instancia}{instancia de examen}{instancias de examen}{Copia física
digitalizada de un examen asociada a un estudiante concreto, que incluye el
\acs{pdf} escaneado, sus anotaciones, calificaciones y metadatos}

\newdefinition{modelo}{modelo de examen}{modelos de examen}{Variante estructural
de un examen que define el número de problemas, sus puntuaciones, las zonas de
reconocimiento y el \acs{pdf} en blanco de referencia}

\newdefinition{anotacion}{anotación}{anotaciones}{Registro creado por un corrector
sobre una \dfn{instancia} que puede contener texto enriquecido, trazos de
\textit{stylus} o audio, y que opcionalmente se asocia a un problema concreto}

\newdefinition{rubrica}{rúbrica}{rúbricas}{Conjunto de criterios de evaluación con
puntuación fija asociados a un problema, cuyos \textit{checkmarks} el corrector
marca para calcular automáticamente la calificación}

\newdefinition{corrector}{corrector}{correctores}{Profesor al que se le han asignado
una o más \dfnpl{instancia} o problemas de un examen para su evaluación}

\newdefinition{coordinador}{coordinador}{coordinadores}{Profesor con el rol principal
de gestión dentro de una asignatura, responsable de crear exámenes, asignar
correctores, gestionar revisiones y publicar calificaciones}

\newdefinition{gestor}{gestor}{gestores}{Usuario con permisos de administración de la
plataforma, responsable de la gestión de usuarios, cursos académicos y configuración
global del sistema}

\newdefinition{watcher}{watcher}{watchers}{Proceso demonio que monitoriza una o varias
carpetas \acs{sftp} en busca de nuevos ficheros \acs{pdf} para iniciar su ingesta
automática}

\newdefinition{softdelete}{desactivación lógica}{desactivaciones lógicas}{Mecanismo de
baja de un registro que preserva su existencia en la base de datos mediante un campo
indicador, garantizando la integridad referencial de los datos históricos}

\newdefinition{zonaocr}{zona de reconocimiento}{zonas de reconocimiento}{Región
rectangular definida por coordenadas y página sobre el \acs{pdf} en blanco de un
\dfn{modelo}, que indica al motor \acs{ocr} dónde buscar un atributo concreto}

\newdefinition{regla}{regla de asignación}{reglas de asignación}{Configuración que
define qué \dfnpl{corrector} evalúan qué \dfnpl{instancia} o problemas, basándose en
filtros como grupo, \dfn{modelo} o problema}

\newdefinition{endpoint}{endpoint}{endpoints}{Punto de acceso de una \ac{api} \ac{rest},
identificado por una \ac{http} y una ruta, que acepta peticiones y devuelve respuestas}

\newdefinition{worker}{worker}{workers}{Proceso independiente que ejecuta tareas
asíncronas extraídas de una cola de mensajes}

\newdefinition{broker}{broker}{brokers}{Intermediario de mensajes que gestiona la
comunicación entre productores y consumidores de tareas asíncronas}

\newdefinition{pipeline}{pipeline}{pipelines}{Secuencia encadenada de etapas de
procesamiento en la que la salida de una etapa alimenta la entrada de la siguiente}

\newdefinition{bucket}{bucket}{buckets}{Contenedor lógico de almacenamiento de objetos en un
servicio de almacenamiento, utilizado para organizar y gestionar los archivos del sistema}

\newdefinition{multi-tenant}{multi-tenant}{multi-tenant}{Arquitectura de software en la que una única
instancia de la aplicación sirve a múltiples organizaciones o clientes, manteniendo la
separación de datos y configuraciones entre ellos}

\newdefinition{single-tenant}{single-tenant}{single-tenant}{Arquitectura de software en la que cada
organización o cliente tiene su propia instancia dedicada de la aplicación, con su propia
base de datos y configuración}

\newdefinition{polling}{polling}{polling}{Técnica de comunicación en la que un proceso cliente
realiza peticiones periódicas a un servidor para comprobar si hay nuevos datos o eventos disponibles}
